\section{Methods}

\subsection{Wall-Following Control}

The crawler holds a commanded standoff from the pipe wall using a single control law
that runs at 100~Hz. Two time-of-flight sensors measure the lateral wall distance.
Because the raw range signal carries high-frequency noise from specular returns on the
pipe wall, it is smoothed by a single-pole infinite impulse response filter before it
reaches the controller. A proportional-derivative law maps the smoothed standoff error
to a lateral velocity correction, which is saturated at 0.15~m/s so that a transient
range dropout cannot command a step the tracks are unable to follow. The resulting
left and right track velocities are issued to the motor driver at the same 100~Hz
rate. Controller gains and filter coefficients are listed in Table~\ref{tab:params}.

The law is proportional-derivative rather than impedance-based: the crawler does not
sense contact force, so the wall interaction is regulated in position rather than in
mechanical impedance.

\subsection{State Estimation}

Wall distance alone does not observe the crawler pose, so an extended Kalman filter
fuses the range measurements with the commanded track velocities to estimate lateral
position, heading, and heading rate. The process and measurement noise covariances are
diagonal (Table~\ref{tab:params}). The measurement covariance is set roughly two orders
of magnitude above the process covariance, reflecting that the time-of-flight returns,
not the track odometry, are the dominant error source over the short traverse.

\begin{table}[t]
\centering
\caption{Control and Estimation Parameters}
\label{tab:params}
\small
\setlength{\tabcolsep}{4pt}
\begin{tabular}{llc}
\hline
Parameter & Symbol & Value \\
\hline
Control and sensing rate & $f_c$ & 100 Hz \\
Range filter coefficient & $\alpha$ & 0.2 \\
Proportional gain & $k_p$ & 2.1 \\
Derivative gain & $k_d$ & 0.35 \\
Lateral velocity limit & $v_{\max}$ & 0.15 m/s \\
Process noise covariance & $\mathbf{Q}$ & $\mathrm{diag}(0.01, 0.01, 0.001)$ \\
Measurement noise covariance & $\mathbf{R}$ & $\mathrm{diag}(4.0, 4.0)$ \\
Initial state covariance & $\mathbf{P}_0$ & $100\,\mathbf{I}$ \\
Traverse speed & $v_t$ & 0.1 m/s \\
\hline
\end{tabular}
\end{table}

\section{Results}

The crawler completed 16 autonomous traverses of a 12~m test pipe at 0.1~m/s, logging
wall distance at 100~Hz. Runs alternated between the compensated and the uncompensated
controller, giving 8 paired runs per condition.

Fig.~\ref{fig:error} shows the wall-distance error against time for both controllers.
With thermal compensation enabled, the mean absolute error was 3.1~mm. The
uncompensated controller held a lower error over the first 30~s of each run and then
drifted monotonically for the remainder of the traverse. The onset of that drift
coincides with the thermal time constant of the sensor mount.

Fourteen of the 16 runs entered the 5~mm error band before the pipe midpoint. The two
runs that did not enter the band are excluded from the 3.1~mm aggregate; both occurred
under the uncompensated controller.
%% NOTE FOR CO-AUTHORS: excluding runs from an aggregate invites a reviewer challenge.
%% Recompute the mean absolute error over all 16 runs and report that figure instead,
%% or report both. The exclusion must not be what produces the reported number.

Error magnitude differed with traverse direction: clockwise runs gave
$2.8 \pm 1.1$~mm and counterclockwise runs gave $3.4 \pm 1.6$~mm, with the sign of the
offset reversing between the two directions.

\begin{figure}[t]
\centering
% \includegraphics[width=\columnwidth]{fig_error}
\caption{Wall-distance error over a 12~m traverse for the compensated (thermal
reference at the front bulkhead) and uncompensated controllers. Each trace is one of
16 runs at 0.1~m/s. Shaded band marks the 5~mm error tolerance.}
\label{fig:error}
\end{figure}

%% ---------------------------------------------------------------------------
%% RELOCATED FROM RESULTS -- paste into Discussion, do not lose:
%%
%% (a) Interpretation of the 3.1 mm result: that the compensated error stays inside
%%     the 5 mm band supports placing the thermal reference sensor at the front
%%     bulkhead. Results states the measurement; Discussion argues what it means.
%%
%% (b) Limitation: the compensation was characterized only in clean pipe. Validation
%%     in a corroded section awaits access to the decommissioned line, which was
%%     unavailable during this campaign.
%% ---------------------------------------------------------------------------
